\documentclass[10pt,a4paper]{amsart}
\usepackage{amsmath,amssymb,amsthm,graphicx}
% For commutative diagrams we use Xy-pic
\usepackage[all]{xy}

% We use 2cell for 2-commutative diagrams.
\xyoption{2cell}
% Theorem environments.
%
\theoremstyle{plain}
\newtheorem{theorem}[subsection]{Theorem}
\newtheorem{proposition}[subsection]{Proposition}
\newtheorem{lemma}[subsection]{Lemma}

\theoremstyle{definition}
\newtheorem{definition}[subsection]{Definition}
\newtheorem{example}[subsection]{Example}
\newtheorem{exercise}[subsection]{Exercise}
\newtheorem{situation}[subsection]{Situation}

\theoremstyle{remark}
\newtheorem{remark}[subsection]{Remark}
\newtheorem{remarks}[subsection]{Remarks}

\numberwithin{equation}{subsection}

% Macros
%
\def\lim{\mathop{\mathrm{lim}}\nolimits}
\def\colim{\mathop{\mathrm{colim}}\nolimits}
\def\Spec{\mathop{\mathrm{Spec}}}
\def\Hom{\mathop{\mathrm{Hom}}\nolimits}
\def\Ext{\mathop{\mathrm{Ext}}\nolimits}
\def\SheafHom{\mathop{\mathcal{H}\!\mathit{om}}\nolimits}
\def\SheafExt{\mathop{\mathcal{E}\!\mathit{xt}}\nolimits}
\def\Sch{\mathit{Sch}}
\def\Mor{\mathop{\mathrm{Mor}}\nolimits}
\def\Ob{\mathop{\mathrm{Ob}}\nolimits}
\def\Sh{\mathop{\mathit{Sh}}\nolimits}
\def\NL{\mathop{N\!L}\nolimits}
\def\CH{\mathop{\mathrm{CH}}\nolimits}
\def\proetale{{pro\text{-}\acute{e}tale}}
\def\etale{{\acute{e}tale}}
\def\QCoh{\mathit{QCoh}}
\def\Ker{\mathop{\mathrm{Ker}}}
\def\Im{\mathop{\mathrm{Im}}}
\def\Coker{\mathop{\mathrm{Coker}}}
\def\Coim{\mathop{\mathrm{Coim}}}

% Boxtimes
%
\DeclareMathSymbol{\boxtimes}{\mathbin}{AMSa}{"02}

%
% Macros for moduli stacks/spaces
%
\def\QCohstack{\mathcal{QC}\!\mathit{oh}}
\def\Cohstack{\mathcal{C}\!\mathit{oh}}
\def\Spacesstack{\mathcal{S}\!\mathit{paces}}
\def\Quotfunctor{\mathrm{Quot}}
\def\Hilbfunctor{\mathrm{Hilb}}
\def\Curvesstack{\mathcal{C}\!\mathit{urves}}
\def\Polarizedstack{\mathcal{P}\!\mathit{olarized}}
\def\Complexesstack{\mathcal{C}\!\mathit{omplexes}}
% \Pic is the operator that assigns to X its picard group, usage \Pic(X)
% \Picardstack_{X/B} denotes the Picard stack of X over B
% \Picardfunctor_{X/B} denotes the Picard functor of X over B
\def\Pic{\mathop{\mathrm{Pic}}\nolimits}
\def\Picardstack{\mathcal{P}\!\mathit{ic}}
\def\Picardfunctor{\mathrm{Pic}}
\def\Deformationcategory{\mathcal{D}\!\mathit{ef}}
\usepackage{cleveref}
\crefname{originaldefinition}{original definition}{original definitions}
\Crefname{originaldefinition}{Original definition}{Original definitions}
\makeatletter
\newlabel{situation-setup-base-change}{{0FVG}{source}{}{sourcepointer-1}{}}
\newlabel{situation-setup-base-change@cref}{{[originaldefinition][0][]0FVG}{[1][0][]source}{}{}{}}
\makeatother
\crefname{originallemma}{original lemma}{original lemmas}
\Crefname{originallemma}{Original lemma}{Original lemmas}
\makeatletter
\newlabel{lemma-pullback-coherent}{{02RE}{source}{}{sourcepointer-2}{}}
\newlabel{lemma-pullback-coherent@cref}{{[originallemma][0][]02RE}{[1][0][]source}{}{}{}}
\makeatother
\crefname{originallemma}{original lemma}{original lemmas}
\Crefname{originallemma}{Original lemma}{Original lemmas}
\makeatletter
\newlabel{lemma-dimension-base-change}{{0FVH}{source}{}{sourcepointer-3}{}}
\newlabel{lemma-dimension-base-change@cref}{{[originallemma][0][]0FVH}{[1][0][]source}{}{}{}}
\makeatother
\crefname{originallemma}{original lemma}{original lemmas}
\Crefname{originallemma}{Original lemma}{Original lemmas}
\makeatletter
\newlabel{lemma-cycle-closed-coherent}{{02QY}{source}{}{sourcepointer-4}{}}
\newlabel{lemma-cycle-closed-coherent@cref}{{[originallemma][0][]02QY}{[1][0][]source}{}{}{}}
\makeatother
\crefname{originallemma}{original lemma}{original lemmas}
\Crefname{originallemma}{Original lemma}{Original lemmas}
\makeatletter
\newlabel{coherent-lemma-coherent-Noetherian}{{01XZ}{source}{}{sourcepointer-5}{}}
\newlabel{coherent-lemma-coherent-Noetherian@cref}{{[originallemma][0][]01XZ}{[1][0][]source}{}{}{}}
\makeatother
\crefname{originallemma}{original lemma}{original lemmas}
\Crefname{originallemma}{Original lemma}{Original lemmas}
\makeatletter
\newlabel{modules-lemma-pullback-finite-type}{{01B6}{source}{}{sourcepointer-6}{}}
\newlabel{modules-lemma-pullback-finite-type@cref}{{[originallemma][0][]01B6}{[1][0][]source}{}{}{}}
\makeatother
\crefname{originallemma}{original lemma}{original lemmas}
\Crefname{originallemma}{Original lemma}{Original lemmas}
\makeatletter
\newlabel{lemma-length-finite}{{02QW}{source}{}{sourcepointer-7}{}}
\newlabel{lemma-length-finite@cref}{{[originallemma][0][]02QW}{[1][0][]source}{}{}{}}
\makeatother
\crefname{originallemma}{original lemma}{original lemmas}
\Crefname{originallemma}{Original lemma}{Original lemmas}
\makeatletter
\newlabel{algebra-lemma-pullback-module}{{02M1}{source}{}{sourcepointer-8}{}}
\newlabel{algebra-lemma-pullback-module@cref}{{[originallemma][0][]02M1}{[1][0][]source}{}{}{}}
\makeatother
\title{Base change of cycles and of the k-cycle of a coherent sheaf, with the expected dimension shift}
\author{The Stacks Project Authors}
\date{stacks-project:chow.tex:local-snapshot-20261012; GNU FDL 1.2 or later}
\begin{document}
\maketitle
\noindent\textit{Editorial note.} All delivered mathematical text is the author's own, copied from the pinned Stacks Project chapter file chow.tex; no statement or proof step was rewritten, reordered or removed.\par\medskip
\noindent\textit{Editorial note.} Preamble provenance: the pinned chapter file opens with an input of the shared preamble and carries no preamble text of its own, so the project's single Stacks-wide preamble staged at sources/stacks-chow/preamble.tex is excerpted verbatim as two spans, lines 14-18 (Xy-pic setup) and lines 154-217 (the native theorem declarations and mathematical macros). Nothing was added, removed or redefined; the wrapper class and the amsmath/amssymb/amsthm/graphicx base packages are supplied by the assembler.\par\medskip
\noindent\textit{Editorial note.} The author opens the delivered proof with a sentence kept verbatim in the text (source lines 15211-15213): "The proof is exactly the same is the proof of Lemma \textbackslash{}ref\{lemma-pullback-coherent\} and we suggest the reader skip it." The argument itself is then written out, namely the dimension input, the reduction of part (1) to part (2) and the full coefficient computation for part (2).\par\medskip
\noindent\textit{Editorial note.} Canonical location of the counted proposition: Stacks tag 0FVI (sources/stacks-chow/pinned-original-source.tex line 15193).\par\medskip
\noindent\textit{Allowed standard prerequisites (editorial).}\begin{itemize}
\item The base change situation for the cycle theory and the associated dimension function delta'. Same chapter, Situation labelled situation-setup-base-change (source line 15123), invoked in the statement. Fixes the constant c, the morphism S' -> S, the base change X' -> S' and the function delta' with respect to which all dimensions are measured.
\item Dimension of a base change and passage between a closed subscheme and its structure sheaf. Same chapter, Lemma labelled lemma-dimension-base-change (source line 15139) and Lemma labelled lemma-cycle-closed-coherent (source line 1776), invoked at proof lines 15216-15219. Supply the dimension bound dim\_\{delta'\}(Z') <= k + c and the derivation of the closed subscheme statement from the coherent sheaf statement.
\item Length of a pullback module along a flat local map. Algebra, Lemma labelled algebra-lemma-pullback-module (Stacks tag 02M1), invoked at proof line 15251. Supplies the equality length\_B(B tensor\_A M) = length\_A(M) length\_B(B/m\_AB) whose two sides are the two coefficients computed in the counted span.
\item Local Noetherianity of X and X' and coherence of the base change. Cohomology of Schemes, Lemma labelled coherent-lemma-coherent-Noetherian (Stacks tag 01XZ) and Modules, Lemma labelled modules-lemma-pullback-finite-type (Stacks tag 01B6), invoked at proof lines 15223-15228. Supply that F' is coherent so that the cycle [F']\_\{k+c\} makes sense.
\end{itemize}\medskip
\ifdefined\hypertarget\hypertarget{sourcepointer-1}{}\fi
\noindent\textit{Source reference.} Original definition 0FVG: sources/stacks-chow/pinned-original-source.tex, source lines 15123--15123. Same-chapter pointer into the author environment situation at source line 15123; invoked at proof line None of the counted proof. It is retained only as a navigation pointer; every load-bearing step is inside the delivered span.\par\medskip
\ifdefined\hypertarget\hypertarget{sourcepointer-2}{}\fi
\noindent\textit{Source reference.} Original lemma 02RE: sources/stacks-chow/pinned-original-source.tex, source lines 2331--2331. Same-chapter pointer into the author environment lemma at source line 2331; invoked at proof line 15212 of the counted proof. It is retained only as a navigation pointer; every load-bearing step is inside the delivered span.\par\medskip
\ifdefined\hypertarget\hypertarget{sourcepointer-3}{}\fi
\noindent\textit{Source reference.} Original lemma 0FVH: sources/stacks-chow/pinned-original-source.tex, source lines 15139--15139. Same-chapter pointer into the author environment lemma at source line 15139; invoked at proof line 15217 of the counted proof. It is retained only as a navigation pointer; every load-bearing step is inside the delivered span.\par\medskip
\ifdefined\hypertarget\hypertarget{sourcepointer-4}{}\fi
\noindent\textit{Source reference.} Original lemma 02QY: sources/stacks-chow/pinned-original-source.tex, source lines 1776--1776. Same-chapter pointer into the author environment lemma at source line 1776; invoked at proof line 15218 of the counted proof. It is retained only as a navigation pointer; every load-bearing step is inside the delivered span.\par\medskip
\ifdefined\hypertarget\hypertarget{sourcepointer-5}{}\fi
\noindent\textit{Source reference.} Original lemma 01XZ: sources/stacks-chow/tags.txt, source lines 2538--2538. Cross-chapter standard prerequisite located in the project tag list staged in the local snapshot; invoked at proof line 15224 of the counted proof. It is retained only as a navigation pointer; every load-bearing step is inside the delivered span.\par\medskip
\ifdefined\hypertarget\hypertarget{sourcepointer-6}{}\fi
\noindent\textit{Source reference.} Original lemma 01B6: sources/stacks-chow/tags.txt, source lines 1715--1715. Cross-chapter standard prerequisite located in the project tag list staged in the local snapshot; invoked at proof line 15226 of the counted proof. It is retained only as a navigation pointer; every load-bearing step is inside the delivered span.\par\medskip
\ifdefined\hypertarget\hypertarget{sourcepointer-7}{}\fi
\noindent\textit{Source reference.} Original lemma 02QW: sources/stacks-chow/pinned-original-source.tex, source lines 1702--1702. Same-chapter pointer into the author environment lemma at source line 1702; invoked at proof line 15240 of the counted proof. It is retained only as a navigation pointer; every load-bearing step is inside the delivered span.\par\medskip
\ifdefined\hypertarget\hypertarget{sourcepointer-8}{}\fi
\noindent\textit{Source reference.} Original lemma 02M1: sources/stacks-chow/tags.txt, source lines 3428--3428. Cross-chapter standard prerequisite located in the project tag list staged in the local snapshot; invoked at proof line 15251 of the counted proof. It is retained only as a navigation pointer; every load-bearing step is inside the delivered span.\par\medskip
\begin{lemma}
\label{lemma-pullback-coherent-base-change}
In Situation \ref{situation-setup-base-change} let $X \to S$
locally of finite type and let $X' \to S$ be the base change by $S' \to S$.
\begin{enumerate}
\item Let $Z \subset X$ be a closed subscheme with
$\dim_\delta(Z) \leq k$ and base change $Z' \subset X'$. Then we have
$\dim_{\delta'}(Z')) \leq k + c$
and $[Z']_{k + c} = g^*[Z]_k$ in $Z_{k + c}(X')$.
\item Let $\mathcal{F}$ be a coherent sheaf on $X$ with
$\dim_\delta(\text{Supp}(\mathcal{F})) \leq k$ and base
change $\mathcal{F}'$ on $X'$.
Then we have $\dim_\delta(\text{Supp}(\mathcal{F}')) \leq k + c$
and $g^*[\mathcal{F}]_k = [\mathcal{F}']_{k + c}$
in $Z_{k + c}(X')$.
\end{enumerate}
\end{lemma}

\begin{proof}
The proof is exactly the same is the proof of
Lemma \ref{lemma-pullback-coherent}
and we suggest the reader skip it.

\medskip\noindent
The statements on dimensions follow from
Lemma \ref{lemma-dimension-base-change}.
Part (1) follows from part (2) by Lemma \ref{lemma-cycle-closed-coherent}
and the fact that the base change of the coherent module $\mathcal{O}_Z$
is $\mathcal{O}_{Z'}$.

\medskip\noindent
Proof of (2). As $X$, $X'$ are locally Noetherian we may apply
Cohomology of Schemes, Lemma \ref{coherent-lemma-coherent-Noetherian} to see
that $\mathcal{F}$ is of finite type, hence $\mathcal{F}'$ is
of finite type (Modules, Lemma \ref{modules-lemma-pullback-finite-type}),
hence $\mathcal{F}'$ is coherent
(Cohomology of Schemes, Lemma \ref{coherent-lemma-coherent-Noetherian} again).
Thus the lemma makes sense. Let $W \subset X$ be an integral closed
subscheme of $\delta$-dimension $k$, and let $W' \subset X'$ be
an integral closed subscheme of $\delta'$-dimension $k + c$ mapping into $W$
under $X' \to X$. We have to show that the coefficient $n$ of
$[W']$ in $g^*[\mathcal{F}]_k$ agrees with the coefficient
$m$ of $[W']$ in $[\mathcal{F}']_{k + c}$. Let $\xi \in W$ and
$\xi' \in W'$ be the generic points. Let
$A = \mathcal{O}_{X, \xi}$, $B = \mathcal{O}_{X', \xi'}$
and set $M = \mathcal{F}_\xi$ as an $A$-module. (Note that
$M$ has finite length by our dimension assumptions, but we
actually do not need to verify this. See
Lemma \ref{lemma-length-finite}.)
We have $\mathcal{F}'_{\xi'} = B \otimes_A M$.
Thus we see that
$$
n = \text{length}_B(B \otimes_A M)
\quad
\text{and}
\quad
m = \text{length}_A(M) \text{length}_B(B/\mathfrak m_AB)
$$
Thus the equality follows from
Algebra, Lemma \ref{algebra-lemma-pullback-module}.
\end{proof}

\end{document}
